\subsection{Finite subgroup descent and the branch-order comparison}

\begin{result}{E.44}{Proved}{The two forgetting maps give the same elliptic parameter}
The index-$24$ branch-order cover in E.37 is intrinsically recognized
from the labelled universal-curve extension.  The explicit
ten-point construction below supplies its named genus-zero comparison
and proves $\mathcal J(A)=A$ for every
$A\in\widehat{\mathrm{GT}}$.
Thus the morphism (E.134), including its specified five-point source,
is compatible with the action in E.34--E.39 in the exact outer
fundamental-group diagram of E.37.

The finite-subgroup theorem used in the proof is
\cite[Theorem A]{ET_BZ2024}, applied to the single finite $2$-group
$C_2^2$.  Both pure conjugacy classes of its labelled sphere action
are retained.  The proof uses the actual forgetting maps and
finite group quotients displayed below.
\end{result}
\begin{proof}
\par\medskip\noindent\textbf{The intrinsic four-element set.}\quad
For a torus with two ordered points, the map
$\iota(z)=p_1+p_2-z$ is independent of the choice of origin.
It naturally interchanges the two labels.  Accordingly the
mapping class group allowing that interchange is
$\Gamma_{1,2}\times\langle\iota\rangle$.
The pure completed group is center-free by E.10, so this is also
the unique central lift of the transposition after completion.
In the universal-curve extension, the preimage of this central
order-two group is
\[
 D_4=C_2(x_1)\amalg C_2(x_2)\amalg C_2(x_3)\amalg C_2(x_4),
 \qquad K=\ker(D_4\to C_2),\quad x_i\longmapsto1.
                                                               \tag{E.210}
\]
This is the actual extension of E.9: the orbifold quotient is a
sphere with four order-two points and one puncture, and its
index-two cover is the twice-punctured torus.  Equivalently use
$(D_4\rtimes B_4)/\langle(c,\omega_4)\rangle$ from E.9 with both
parities allowed; its map to
$(B_4/\langle\omega_4\rangle)\times C_2$ has the displayed preimage.

There are exactly four $K$-conjugacy classes of order-two lifts,
represented by $x_1,\ldots,x_4$.  Discretely this is the reduced-word
torsion theorem for a free product.  After completion Theorem A,
for this virtually free group and $p=2$, gives the same assertion.
The classes are distinct in abelianization.  Each $D_4$-class is
one $K$-class, since its centralizer contains its odd generator
and $K\langle x_i\rangle=D_4$.
The $B_4$ action on these classes is its strand permutation.
Ordering them therefore gives exactly
$\widehat\Gamma_{1,2}\to\mathfrak S_4$, with kernel
$\widehat P_4/\overline{\langle\omega_4\rangle}$.
The double-cover construction of E.37 identifies that kernel
with its named five-point group.  The parameter comparison still
requires a map; the remainder of the proof supplies it.

\par\medskip\noindent\textbf{Two preliminary symmetry checks.}\quad
The quotient of the four-point line by $z\mapsto1-z$ has
\[
 U=4z(1-z),\qquad
 [\overline M_{0,4}/C_2]_{U=0}\simeq *,
 \qquad [\overline M_{0,4}/C_2]_{U=\infty}\simeq BC_2 .
                                                               \tag{E.211}
\]
The first boundary is the free orbit $\{0,1\}$ and the second
is the fixed point $\infty$.  Thus the open cusp swap
$U\mapsto1/U$ does not extend to an equivalence of this particular
compactified quotient exit category: its endpoint groups would
identify $1$ with $C_2$.
For the involution family $\{0,\pm1,\pm z\}$, adding the other
fixed pole and exchanging the poles sends $t=z^2$ to $1/t$.
Its quotient parameter $(t+1)^2/(4t)$ stays fixed, so this gives
the identity outer map on that quotient, not its cusp swap.
These calculations only concern these two specified candidates.

\par\medskip\noindent\textbf{The ten-point family and the actual centralizer.}\quad
Use the labelled configuration
\[
\begin{gathered}
 H=\langle a,b\rangle\cong C_2^2,\qquad
 a(w)=-w,\quad b(w)=1/w,\quad c=ab,\\
 A=(0,\infty),\quad B=(1,-1),\quad C=(i,-i),\qquad
 D=(z,-z,1/z,-1/z),\\
 z\in\mathbb P^1\setminus\{0,\infty,1,-1,i,-i\}.
\end{gathered}                                                   \tag{E.212}
\]
The pairs $A,B,C$ are the fixed pairs of $a,b,c$; $D$ is a regular
$H$-orbit.  Let $G$ be the inverse image of
$C_{\mathfrak S_{10}}(H)$ in the full sphere mapping class group,
and set $N=C_G(H)$.  This permutation centralizer has order
$2^3\cdot4=32$: it allows each fixed-pair interchange and an
$H$-equivariant translation of $D$, while preserving the names
of $A,B,C$.

Forgetting the six fixed labels gives
\[
 f_4:N\xrightarrow{\ \sim\ }
 C_{\Gamma_{0,4}^{\rm full}}(H_D)
      =\Gamma_{0,4}\times H_D,\qquad
 \bar f_4:N/H\xrightarrow{\ \sim\ }\Gamma_{0,4}.
                                                               \tag{E.213}
\]
Here $H_D$ is the canonical normal Klein-four group of double
transpositions of the remaining quartet.  We first justify the
mapping-class assertion, rather than infer it from coarse moduli.

Remove all ten source punctures and all four quotient punctures.
The actual unramified cover gives
$1\to\Pi\to E\to H\to1$, with $\Pi=F_9$ and $E=F_3$.
Three quotient peripheral generators $u_A,u_B,u_C$ map to $a,b,c$,
and the fourth maps to $1$.  The center of $\Pi$ and its
centralizer in $E$ are trivial.  Conjugation identifies $E$ with
the inverse image of $H\subset\operatorname{Out}(\Pi)$ in
$\operatorname{Aut}(\Pi)$, exactly as in E.8.
Hence any representative of a mapping class centralizing $H$
normalizes this actual $E$ and induces an automorphism of it
acting identically on $H$.

At a branch cusp $u_A^2$ is an upstairs peripheral element.
Preservation of the named upstairs pair makes its image
$E$-conjugate to the same positive square; the conjugator need
not lie in $\Pi$.  Unique roots in the discrete free group $E$
make the image of $u_A$ conjugate to $u_A$.  The same argument
applies to $B,C$, and the fourth cusp requires no square.
The oriented punctured-sphere Dehn--Nielsen--Baer identification
therefore yields a pure four-point mapping class downstairs.
Its kernel is exactly the deck group $H$, since an inner
downstairs representative restricts to that deck class upstairs.
Conversely a pure downstairs mapping class fixes the homomorphism
$E\to H$: all peripheral classes are preserved and $H$ is
abelian.  It lifts, and its restriction to $\Pi$ is oriented
peripheral, hence is an upstairs mapping class centralizing $H$.
This proves the exact centralizer sequence with actual maps.

The quotient parameter is $u=(z^2+1)^2/(4z^2)$, and the
cross-ratio of the ordered generic quartet is
$\lambda=((z^2-1)/(z^2+1))^2=1-1/u$.
Thus forgetting the six fixed points induces an isomorphism
on the downstairs four-point group.  On the deck kernel it is
the isomorphism $H\to H_D$, proving (E.213).
Its inverse on families is explicit: take the canonical
Klein-four group of a four-point sphere and add its three
fixed pairs, named by their nonidentity elements.
These six points are distinct and avoid the original quartet.
Restriction and this operation are inverse on equivariant
isomorphisms.

The topology is also exact.  Forgetting is defined on all of
the finite-index mixed group $G$ and lands in
$\Gamma_{0,4}\times H_D$.  Every finite quotient of that target
pulls back to a finite quotient of $G$, whose restriction to $N$
is the prescribed quotient under (E.213).
Thus $N$ has its full induced profinite topology.  In fact
forgetting followed by the inverse of (E.213) is a retraction
$G\to N$.  Theorem A gives
$C_{\widehat G}(H)=\overline N=\widehat N$.

\par\medskip\noindent\textbf{The two pure classes and their common correction.}\quad
Theorem A applies to this $G$ with $p=2$.
Pure genus-zero mapping class groups are iterated finitely
generated free-group extensions by the forgetting sequences,
with full kernel topology from characteristic finite quotients.
The Lyndon--Hochschild--Serre comparison gives goodness
inductively; finite classifying models give finite cohomological
type and finite virtual cohomological dimension.
These properties pass to the finite extension $G$, which is of
type $\mathrm{FP}_\infty$.  Discrete group cohomology therefore
commutes with filtered unions of modules; continuous profinite
cochains commute with these unions by compactness.  Every discrete
continuous $\mathbb F_2$-module is a union of finite submodules,
since the orbit of each of its vectors is finite.
The required comparison for arbitrary discrete
$\mathbb F_2[[\widehat G]]$-modules follows by filtered union
of finite submodules and the same finite-type property.
Thus Theorem A(i) geometrizes the whole finite $2$-subgroup $H$,
and (iii) gives its centralizer closure.
These are the prime-group assertions; no nonsolvable step is used.

There are exactly two pure conjugacy classes for the elementwise
label permutations in (E.212).
Finite realization on the punctured sphere is
\cite[Theorem 1 and its punctured extension, p.~236]{ET_Kerckhoff1983}.
Uniformize the realized sphere and normalize
$A_1=0,A_2=\infty,B_1=1$.
The commuting involutions become $a(w)=-w,b(w)=1/w$.
Consequently $B_2=-1$ and
\[
 (C_1,C_2)=(i,-i)\ \text{or}\ (-i,i),\qquad
 P=(C_1\ C_2).
                                                               \tag{E.214}
\]
The generic orbit is determined by its first labelled point $z$.
Its allowed domain in (E.212) is connected; a path lifts to an
equivariant isotopy carrying its four labels.  Each sign is
therefore one pure class, and distinct signs cannot meet while
the ten labels remain distinct.  Relabeling by $P$ exchanges
the two classes.

Let $F$ be the genus-zero tower action.
It acts identically on the finite label-permutation quotient.
Apply Theorem A to $F(H)$ inside $\widehat G$.
Replace the permutation of the resulting conjugator by a
discrete element of $G$; the remaining conjugator is pure.
Because that permutation centralizes $H$, the resulting discrete
subgroup has exactly the elementwise permutation data (E.212).
The preceding classification gives a further pure conjugation
and, when needed, the relabeling $P$.
After this correction $F$ fixes $H$ pointwise, since each of its
elements is determined by its faithful permutation on the labels.
Represent $P$ by a half-twist supported in a disk containing only
$C_1,C_2$.  After forgetting those two labels, that disk is unmarked
and the half-twist is isotopic to the identity relative to every
remaining label.  Thus $P$ is literally killed by (E.213) and by
both forgetting maps below.
The pure corrections only supply their usual inner basepoint
comparisons.  We have therefore retained both pure classes,
rather than assumed that a single one is invariant.

\par\medskip\noindent\textbf{The colored target and the choice subgroups.}\quad
For a quartet with two distinguished two-element blocks put
$C_4=\mathfrak S_2\times\mathfrak S_2$, and let $D_8$ preserve the
unordered partition into those blocks.  Let $V$ be the normal
Klein-four subgroup of the full four-point group.  Then
\[
 R=\Gamma_{0,4}^{D_8}/V\cong C_2*\mathbb Z,\qquad
 \Gamma_{0,4}^{C_4}/\langle\text{double flip}\rangle
                  \xrightarrow{\ \sim\ }R.
                                                               \tag{E.215}
\]
Indeed $C_4\cap V$ is the order-two double flip and $C_4V=D_8$.
Thus the actual colored forgetting image suffices; no
interchange of the two colors is silently added.

The quotient by $V$ is intrinsic also after completion.
A finite normal subgroup intersects the torsion-free pure
$\widehat F_2$ trivially, so centralizes it.
An element centralizing that pure group acts trivially on its
abelianization; the quotient $\mathfrak S_4/V=\mathfrak S_3$
acts faithfully there by permuting the three peripheral vectors.
Hence that element lies in $\widehat F_2 V$, and centerlessness
of $\widehat F_2$ leaves exactly $V$.
Thus $V$ is the maximal finite normal subgroup and is
characteristic for the induced action.

Choose the $b$-orbit $O_b=\{z,1/z\}\subset D$ and let $N_b$
stabilize it.  The action of $H$ on the two such choices gives
\[
 N=HN_b,\quad H\cap N_b=\langle b\rangle,\quad
 N_b/\langle b\rangle\xrightarrow{\ \sim\ }N/H,
 \qquad q_b:N/H\longrightarrow R.
                                                               \tag{E.216}
\]
To define the last map, invert the preceding isomorphism,
retain $A\amalg O_b$ by ordinary forgetting, include the colored
target into the $D_8$ target, and quotient by $V$.
The element $b$ simultaneously flips both retained pairs and
dies in $V$.  Thus the map is well-defined.
Explicitly lift a class to $N_b$ by multiplying by an element
of $H$; two such lifts differ by $b$, which has been killed.
Similarly choose $O_a=\{z,-z\}$, retain $B\amalg O_a$ and kill
its double flip $a$, giving $q_a:N/H\to R$ with the same named
target partition.
Both choice subgroups have index two, and all quotient kernels
just used are finite.  These actual forgetting constructions
therefore give the same continuous maps on full completions.
The correction $P$ affects neither chosen orbit, neither
retained pair, nor the generic quartet.

\par\medskip\noindent\textbf{The two rational maps and surjectivity.}\quad
For two unordered pairs choose a cross-ratio $\ell$ such that
an interchange in either pair sends it to $1/\ell$.
The effective quotient coordinate is $(\ell+1)^2/(4\ell)$.
In our two cases
\[
\begin{aligned}
 \ell_b&=z^2,& U_b&=\frac{(z^2+1)^2}{4z^2}=u,\\
 \ell_a&=\left(\frac{z-1}{z+1}\right)^2,&
 U_a&=\frac{(z^2+1)^2}{(z^2-1)^2}=\frac{u}{u-1},\\
 S(U)&=\frac{U}{U-1},& q_a&=S_Rq_b=q_bS_X.
\end{aligned}                                                   \tag{E.217}
\]
Here $R$ has a cone of order two at $U=0$ and cusps at $1,\infty$.
The source is $X=\mathbb P^1_u\setminus\{0,1,\infty\}$, identified
with the named generic quartet by $\lambda=1-1/u$.
Consequently $q_b$ is the inclusion of the complement of the
orbifold point into $R$ and is onto on fundamental groups.
For positive cusp loops $v$ at $\infty$ and $d$ at $1$, its
presentation is
$F(v,d)\twoheadrightarrow\langle v,d\mid(vd)^2=1\rangle$.
The kernel is the normal closure of $(vd)^2$, and after
completion its closed normal closure.
The cusp swap $S_R$ has outer representative $v\leftrightarrow d$.

The named pure index-two subgroup in the colored target is
$K_{\rm std}=\ker(v\mapsto0,d\mapsto1)=\widehat F(v,d^2)$:
the cusp $1$ is ramified and $\infty$ has two lifts.
This is exactly the reading in E.28.
Thus the induced action on $R$ is $\iota(A)$.
For uniqueness, align two representatives on $K_{\rm std}$ by
its inner basepoint comparison; their quotient fixes that
normal subgroup pointwise and is therefore the identity
because its centralizer is trivial, as proved in E.27.

\par\medskip\noindent\textbf{Naturality and equality of parameters.}\quad
After (E.214), the tower equivalence preserves $H$, its full
completed centralizer, both finite choice subgroups, the
retained labels, the actual forgetting maps and the finite
normal quotient $V$.  Hence the two outer-commutative diagrams
are
\[
 q_b\alpha_A=\beta_Aq_b,\qquad
 q_a\alpha_A=\beta_Aq_a,\qquad
 [\alpha_A]=A,\quad[\beta_A]=\iota(A).
                                                               \tag{E.218}
\]
Each equality permits its one target inner comparison, coming
from the corrected source conjugation and that forgetting map.
The action on $N/H$ is the named $A$ by (E.213), so the source
action in the two diagrams is the same.
Use $q_a=S_Rq_b$ in the two equations.
They give $\beta_AS_Rq_b=S_R\beta_Aq_b$ up to one target inner
automorphism.  Surjectivity of $q_b$ gives the equality on all
of $R$, so $[\beta_A]$ commutes with $[S_R]$.
The injective comparison of E.28 yields
\[
 \iota(\mathcal J(A))=S_R\iota(A)S_R=\iota(A),
                  \qquad \mathcal J(A)=A.
                                                               \tag{E.219}
\]
E.37's already proved equivalence now supplies precisely the
named branch-order cover compatibility.

The executed program \texttt{b\_branch\_descent.py} verifies the
ten-point permutations, finite choices, $C_4V=D_8$ and its
intersection, and all six rational identities.
It also computes the fixed-pair discriminants
$4\lambda$, $4(1-\lambda)$, $4\lambda(\lambda-1)$.
Their monodromies around $0,1$ flip $(A,C)$ and $(B,C)$,
respectively, giving the even-flip image and the two classes
in (E.214).
There is also an exact permutation check of the pure-class
count.  Under (E.213), the two free generators of the pure
four-point group give those two independent even flips.
Together with the four deck translations on $D$, this makes
the permutation image of $N$ have order $16$ inside the
order-$32$ group $C_{\mathfrak S_{10}}(H)$.
Within $G$, every normalizer of $H$ centralizes it, since
its conjugation fixes each faithful label permutation of $H$.
The pure-conjugacy double cosets therefore number $32/16=2$;
the transposition $P$ represents the other coset.
The group, topology and naturality proofs above do not rely
on inferring an infinite-group assertion from those finite checks.
\end{proof}

